% Section 12.1, from lectures/L12/appendix/a_review.md. Each lecture the review points back to is
% the chapter of the same number, and each "L04 A.3" is the section it became (\secref{c4:a3}).

\section{Repetition}
\label{c12:app:a}

\subsection{Så används det här avsnittet}
\label{c12:a1}
Kursen har gått fort: fyra ämnen på elva kapitel. Det här avsnittet går igenom dem igen, ordnat
efter kursplanens lärandemål, så att du kan se vad du kan och vad du behöver läsa om.

Varje avsnitt har tre delar:
\begin{itemize}
  \item \textbf{Det viktigaste}, kort, med de metoder som provet bygger på.
  \item \textbf{Kan du \dots?}, en lista att bocka av. Svara på varje fråga på papper, utan att
    titta, och kontrollera sedan mot avsnittet som hänvisningen pekar på.
  \item \textbf{Vanliga fel}, de misstag som oftast kostar poäng.
\end{itemize}

Provets form, poäng och betygsgränser står i bilaga~\ref{app:exam}. Där finns också ett övningsprov,
i bilaga~\ref{app:practice}, med lösningsförslag i bilaga~\ref{app:practice:answers}. Gör det under
två timmar, utan miniräknare, med referensbladet i bilaga~\ref{app:avr} bredvid, precis som på det
riktiga provet.

\subsection{Talsystem och binära koder}
\label{c12:a2}
\emph{Lärandemål: beskriva binära, hexadecimala och decimala talsystem, och omvandla talvärden
mellan dem.} (Kapitel~\ref{c1:ch}.)

\subsubsection{Det viktigaste}
\begin{itemize}
  \item I ett \textbf{positionssystem} är varje siffra värd siffran gånger talbasen upphöjd till
    positionen, räknat från 0 längst till höger. Binärt har basen 2, hexadecimalt 16.
  \item \textbf{Binärt till decimalt}: addera vikterna för ettorna. \code{1010 1110} = 128 + 32 + 8
    + 4 + 2 = 174.
  \item \textbf{Decimalt till binärt}: dra av den största vikt som får plats, om och om igen, eller
    dela med 2 och läs resterna nedifrån (\secref{c1:a4}).
  \item \textbf{Binärt och hexadecimalt}: fyra bitar är en hexadecimal siffra. \code{1011 0110} =
    \code{0xB6}. Omvandla alltid mellan decimalt och hexadecimalt \textbf{via binärt} om du är
    osäker (\secref{c1:a6}).
  \item \textbf{$n$ bitar} räcker till \code{2^n} värden, 0 till \code{2^n - 1}: en byte till
    0--255.
  \item \textbf{Binär addition} görs som för hand, med minnessiffra. En summa som inte får plats i
    åtta bitar lämnar en minnessiffra, carry (\secref{c1:a7}).
  \item \textbf{Koder}: BCD skriver varje decimal siffra med fyra bitar, Graykod ändrar en bit i
    taget, och ASCII ger varje tecken ett tal: \code{'0'} = \code{0x30}, \code{'A'} = \code{0x41}
    (\secref{c1:app:b}).
\end{itemize}

\subsubsection{Kan du \dots?}
\begin{itemize}
  \item \dots\ omvandla 1100 1011, 142 och \code{0x5E} till de två andra talsystemen, utan
    miniräknare?
  \item \dots\ förklara varför en hexadecimal siffra motsvarar exakt fyra bitar?
  \item \dots\ addera två 8-bitars tal binärt och säga om det blir en minnessiffra?
  \item \dots\ skriva 59 i BCD, och säga varför BCD inte är samma sak som 59 binärt?
  \item \dots\ ange ASCII-koderna för tecknen \code{'4'} och \code{'C'}?
\end{itemize}

\subsubsection{Vanliga fel}
\begin{itemize}
  \item Att räkna positionerna från 1 i stället för från 0, så att varje vikt blir dubbelt för
    stor.
  \item Att gruppera bitar i fyror från vänster. Grupperna räknas \textbf{från höger}: \code{1 0110}
    är \code{0x16}.
  \item Att tro att BCD för 59 är 59 binärt. BCD är \code{0101 1001}, binärt är \code{0011 1011}.
\end{itemize}

\subsection{Grindar, boolesk algebra och kontaktnät}
\label{c12:a3}
\emph{Lärandemål: förklara grindar och logiska funktioner; beskriva logisk algebra och förenklingar
av booleska uttryck.} (Kapitel~\ref{c2:ch} och~\ref{c3:ch}.)

\subsubsection{Det viktigaste}
\begin{itemize}
  \item Sju grindar: NOT, AND, OR, NAND, NOR, XOR och XNOR. NAND, NOR och XNOR är inverserna av AND,
    OR och XOR, och XOR är 1 när ingångarna är olika.
  \item \textbf{Notation}: \code{A'} för NOT, \code{AB} för AND, \code{A + B} för OR. NOT binder
    hårdast, sedan AND, sist OR.
  \item \textbf{Från sanningstabell till uttryck}: en AND-term per etta, med varje variabel direkt
    eller inverterad; OR:a ihop termerna. Det ger ett uttryck på \textbf{SP-form}
    (\secref{c2:a6}).
  \item \textbf{Förenkla} genom att bryta ut det gemensamma och använda \code{A + A' = 1}. Tre lagar
    räcker långt: \code{A + AB = A}, \code{A + A'B = A + B} och De Morgan (\secref{c2:a4}).
  \item \textbf{De Morgan}: \code{(A + B)' = A'B'} och \code{(AB)' = A' + B'}. Bryt strecket, byt
    tecknet (\secref{c2:a5}).
  \item \textbf{Analys av ett grindnät}: skriv uttrycket vid varje grinds utgång, från ingångarna
    mot utgången, och räkna sedan ut sanningstabellen rad för rad (\secref{c2:a8}).
  \item \textbf{Kontaktnät}: kontakter i serie är AND, parallellt OR, och en brytande kontakt är
    NOT. Lampan lyser när det finns en sluten väg från +24~V till 0~V (\secref{c2:app:b} och
    \secref{c3:app:a}).
\end{itemize}

\subsubsection{Kan du \dots?}
\begin{itemize}
  \item \dots\ rita symbolen och skriva sanningstabellen för var och en av de sju grindarna?
  \item \dots\ läsa av ett uttryck på SP-form ur en sanningstabell med tre variabler?
  \item \dots\ förenkla \code{AB + AB' + A'B} till två variabler, och säga vilken lag varje steg
    använder?
  \item \dots\ skriva \code{(A + B)'} utan parentes, och bygga OR av enbart NAND-grindar?
  \item \dots\ skriva uttrycket och sanningstabellen för ett grindnät eller ett kontaktnät du inte
    sett förut?
  \item \dots\ rita ett kontaktnät för \code{S1(S2 + S3')}?
\end{itemize}

\subsubsection{Vanliga fel}
\begin{itemize}
  \item Att glömma att NOT binder hårdare än AND: \code{AB'} är \code{A(B')}, inte \code{(AB)'}.
  \item Att använda De Morgan halvt: invertera variablerna men glömma att byta AND mot OR.
  \item Att tro att en brytande kontakt är öppen: den är \textbf{sluten} tills knappen trycks.
\end{itemize}

\subsection{Karnaughdiagram och IC-kretsar}
\label{c12:a4}
\emph{Lärandemål: använda Karnaughdiagram för analys och minimering av booleska uttryck;
syntetisera digitala kombinatoriska nät med IC-kretsar och kontakter; analysera digitala
kombinatoriska nät.} (Kapitel~\ref{c4:ch} och~\ref{c5:ch}.)

\subsubsection{Det viktigaste}
\begin{itemize}
  \item Ett \textbf{Karnaughdiagram} är sanningstabellen ritad så att grannrutor skiljer sig i
    exakt en variabel. Raderna och kolumnerna står i \textbf{Graykod}: 00, 01, 11, 10.
  \item \textbf{Gruppera} ettorna i rektanglar om 1, 2, 4 eller 8 rutor, så stora som möjligt och så
    få som möjligt. Grupper får överlappa och gå runt kanterna; de fyra hörnen är grannar
    (avsnitt~\ref{c4:a3} och~\ref{c4:a5}).
  \item \textbf{Läs av} varje grupp som de variabler som är konstanta inom den. En grupp om fyra i
    ett diagram med fyra variabler ger en term med två variabler.
  \item \textbf{Don't care}, X, får räknas som 1 om det ger en större grupp, annars som 0.
  \item \textbf{NAND-nät}: ett uttryck på SP-form blir två nivåer NAND-grindar, eftersom
    \code{AB + CD} = \code{((AB)'(CD)')'}.
  \item \textbf{74-serien}: fyra grindar per kapsel (sex inverterare i 74HC04), VCC på ben~14 och
    GND på ben~7. Oanvända ingångar får aldrig flyta (\secref{c4:app:b}).
  \item \textbf{Felsökning}: matningen på varje krets först, sedan ingångarna, sedan signalen grind
    för grind (\secref{c5:app:a}).
\end{itemize}

\subsubsection{Kan du \dots?}
\begin{itemize}
  \item \dots\ rita ett tomt Karnaughdiagram för fyra variabler, med Graykoden rätt på båda
    axlarna?
  \item \dots\ fylla i det från en lista med mintermer, och gruppera det minimalt?
  \item \dots\ förklara varför hörnen i ett diagram med fyra variabler kan bilda en grupp?
  \item \dots\ räkna ut hur många 74HC-kretsar ett nät behöver?
  \item \dots\ förklara vad som händer med en flytande ingång på en 74HC-krets?
\end{itemize}

\subsubsection{Vanliga fel}
\begin{itemize}
  \item Att skriva axlarna i vanlig binär ordning, 00, 01, 10, 11. Då är grannrutorna inte
    grannar.
  \item Att göra en grupp om tre, eller en grupp som inte är en rektangel.
  \item Att behålla en grupp vars alla ettor redan täcks av andra grupper.
\end{itemize}

\subsection{Mikrodatorns uppbyggnad}
\label{c12:a5}
\emph{Lärandemål: redogöra för mikrodatorns uppbyggnad på blockschemanivå.}
(Kapitel~\ref{c6:ch} och~\ref{c7:ch}.)

\subsubsection{Det viktigaste}
\begin{itemize}
  \item En \textbf{mikrodator} består av en \textbf{CPU}, \textbf{minne} och \textbf{in- och
    utportar}, sammankopplade med \textbf{bussar} och drivna av en \textbf{klocka}
    (\secref{c6:b2}).
  \item \textbf{CPU:n} har en \textbf{ALU} som räknar, \textbf{register} som håller värdena den
    räknar med, och en \textbf{styrenhet} som avkodar instruktionerna och styr resten.
  \item \textbf{Programräknaren}, PC, pekar på nästa instruktion. Instruktionscykeln är
    \textbf{hämta, avkoda, utföra}, en gång per instruktion (\secref{c6:b4}).
  \item ATmega328P har \textbf{tre minnen}: programmet i flash, som finns kvar utan ström;
    variablerna och stacken i SRAM, som töms när strömmen bryts; och EEPROM (\secref{c7:a5}).
  \item Varje block är byggt av det kursen redan har gått igenom: ALU:n av grindar, register och PC
    av vippor, minnet av tusentals lås med en adressavkodare (\secref{c6:b10}).
\end{itemize}

\subsubsection{Kan du \dots?}
\begin{itemize}
  \item \dots\ rita mikrodatorns blockschema och säga vad varje block gör?
  \item \dots\ beskriva vad som händer i hämta, avkoda och utföra för en instruktion som
    \code{inc r16}?
  \item \dots\ säga vilket minne programmet ligger i, och vilket stacken ligger i, och vilket av dem
    som finns kvar när strömmen bryts?
  \item \dots\ förklara vad en buss är, och nämna två?
\end{itemize}

\subsubsection{Vanliga fel}
\begin{itemize}
  \item Att blanda ihop registren med dataminnet. Registren är 32 byte inuti CPU:n; SRAM är
    2048 byte utanför den.
  \item Att tro att programmet ligger i SRAM. Det ligger i flash.
\end{itemize}

\subsection{Assemblerprogrammering}
\label{c12:a6}
\emph{Lärandemål: använda assemblerspråk för att programmera en mikrodator.}
(Kapitel~\ref{c7:ch}--\ref{c10:ch}.)

\subsubsection{Det viktigaste}
\begin{itemize}
  \item En rad assembler är en \textbf{etikett}, en \textbf{instruktion} med \textbf{operander}, och
    en \textbf{kommentar}. \textbf{Direktiv} som \code{.include}, \code{.org}, \code{.equ} och
    \code{.def} blir ingen maskinkod.
  \item \code{ldi} laddar en konstant, men bara i \code{r16}--\code{r31}. Tal skrivs decimalt,
    \code{0b...} eller \code{0x...}.
  \item \textbf{Aritmetik och flaggor}: \code{add}, \code{sub}, \code{inc}, \code{dec} och resten
    skriver flaggorna i SREG. \code{C} är minnessiffra eller lån, \code{Z} att resultatet blev noll,
    \code{N} att bit~7 är satt (\secref{c8:a6}). Åtta bitar går runt: 255 + 1 blir 0.
  \item \textbf{Logik}: \code{and}/\code{andi} nollställer bitar, \code{or}/\code{ori} ettställer,
    \code{eor} inverterar, med en mask (\secref{c8:b2}). \code{lsl} multiplicerar med 2 och
    \code{lsr} delar med 2.
  \item \textbf{Hopp}: \code{cp}/\code{cpi} är en subtraktion som bara sparar flaggorna, och det
    villkorliga hoppet läser dem: \code{breq}, \code{brne}, \code{brlo}, \code{brsh}
    (\secref{c9:a4}).
  \item \textbf{Loopar och fördröjningar}: en loop med \code{dec} och \code{brne} tar 3 cykler per
    varv, utom det sista, som tar 2. Vid 16~MHz är en cykel 62,5~ns (\secref{c9:b3}).
  \item \textbf{Portar}: \code{DDRx} väljer riktning, \code{PORTx} sätter en utgång eller pull-upen
    på en ingång, \code{PINx} läser (\secref{c8:c2} och \secref{c10:b2}). En knapp mot jord med
    pull-up läses som 0 när den är nedtryckt.
  \item \textbf{Subrutiner och stacken}: \code{rcall} lägger returadressen på stacken och
    \code{ret} hämtar den. En subrutin sparar varje register den ändrar med \code{push}, och
    återställer dem med \code{pop} i omvänd ordning (\secref{c10:app:c}).
\end{itemize}

\subsubsection{Kan du \dots?}
\begin{itemize}
  \item \dots\ följa ett kort program instruktion för instruktion och ange registren och flaggorna
    \code{C}, \code{Z} och \code{N} efter varje rad?
  \item \dots\ räkna ut hur många cykler en loop tar, och hur lång tid det är vid 16~MHz?
  \item \dots\ skriva ett program som läser en knapp och tänder en lysdiod?
  \item \dots\ skriva en jämförelse och ett villkorligt hopp för \enquote{om \code{r16} är minst
    100}?
  \item \dots\ säga vad \code{SP} är efter ett \code{rcall} och två \code{push}, när den var
    \code{0x08FF} från början?
\end{itemize}

\subsubsection{Vanliga fel}
\begin{itemize}
  \item Att tro att den gula pilen i simulatorn visar instruktionen som just kördes. Den visar
    nästa.
  \item Att lägga en instruktion som ändrar flaggorna mellan \code{cpi} och hoppet.
  \item Att glömma att \code{inc} och \code{dec} inte ändrar \code{C}, och att \code{andi} inte
    heller gör det.
  \item Att räkna en loops sista varv som alla andra. Det sista villkorliga hoppet hoppar inte, och
    tar en cykel mindre.
  \item Att glömma \code{com} när knappar mot jord läses, så att programmet gör precis tvärtom.
\end{itemize}

\subsection{Styrobjekt}
\label{c12:a7}
\emph{Lärandemål: programmera mikrodatorn till att kontrollera styrobjekt.} (Kapitel~\ref{c11:ch}.)

\subsubsection{Det viktigaste}
\begin{itemize}
  \item Ett styrsystem läser \textbf{ingångar}, fattar \textbf{beslut} och sätter
    \textbf{utgångar}. En tillståndstabell och en flödesplan före koden gör programmet enkelt att
    skriva och att kontrollera (\secref{c11:app:b}).
  \item En lysdiod behöver ett \textbf{förkopplingsmotstånd}: \code{(5 V - 2 V) / 220 Ω ≈ 14 mA}.
    Större laster styrs genom en transistor och ett relä (\secref{c11:app:a}).
  \item En knapp \textbf{studsar}; ett program som räknar tryck väntar en kort stund efter varje
    ändring.
\end{itemize}

\subsubsection{Kan du \dots?}
\begin{itemize}
  \item \dots\ räkna ut strömmen genom en lysdiod med ett givet motstånd?
  \item \dots\ rita flödesplanen för ett trafikljus, och översätta den till ett program?
  \item \dots\ förklara varför en knapp kan räknas flera gånger, och hur det förhindras?
\end{itemize}

\subsection{Tips inför provet}
\label{c12:a8}
\begin{itemize}
  \item \textbf{Visa hur du räknar.} Metoden ger poäng även när svaret blir fel på grund av ett
    slarvfel, och ett fel som förs vidare till nästa deluppgift kostar bara en gång.
  \item \textbf{Kontrollera med en sanningstabell} när du har förenklat ett uttryck. Tre variabler
    är åtta rader och tar en minut.
  \item \textbf{Skriv alltid ut talsystemet}: \code{0x2C}, \code{0b0010 1100} eller 44. Ett tal utan
    skrivsätt kan läsas på tre sätt.
  \item \textbf{Följ ett program med en tabell}: en rad per instruktion och en kolumn per register
    och flagga.
  \item \textbf{Läs frågan två gånger.} Frågar den efter ett uttryck, en sanningstabell eller ett
    grindnät? Frågar den efter \code{C} efter en viss instruktion, eller efter hela programmet?
\end{itemize}
